% ============================================
% ОСНОВНАЯ ЧАСТЬ
% Подробное описание выполнения задач этапов 2–4 плана-графика:
%   этап 2 — написание текста ВКР;
%   этап 3 — проверка текста ВКР в системе «Антиплагиат»;
%   этап 4 — подготовка к предзащите, прохождение предзащиты,
%            доработка по итогам предзащиты.
% ============================================

\section{Написание текста ВКР}

\subsection{Выбор инструментов и подготовка шаблона}

Для написания текста выпускной квалификационной работы я использовал
систему вёрстки \LaTeX{}.
Компилировал я работу с помощью XeLaTeX --- это транслятор и диалект издательской системы
\LaTeX{}, работающий на базе движка XeTeX, а сам текст правил в редакторе
Neovim, настроенном под мою работу, на персональном компьютере под
управлением операционной системы Fedora Linux. \LaTeX{} я выбрал по
следующим причинам. Во-первых, моя работа содержит много формул, таблиц,
рисунков и перекрёстных ссылок на них, и в \LaTeX{} все эти элементы
вёрстаются аккуратно и обновляются автоматически при перекомпиляции,
тогда как в Word их пришлось бы выравнивать и перенумеровывать вручную.
Во-вторых, исходный текст работы в \LaTeX{}~--- это обычные текстовые
файлы, поэтому я мог вести их историю в системе контроля версий Git
вместе с исходным кодом самого проекта. В-третьих, разбивка документа на
отдельные файлы позволяла редактировать каждый структурный элемент
изолированно и отслеживать изменения по отдельным разделам.

Готовый шаблон я подготовил заранее, ещё на учебной практике, и на
преддипломной просто взял его за основу. При сборке шаблона я опирался
на нормативные документы~--- ГОСТ~7.32--2017, ГОСТ~Р~7.0.5--2008 и
ГОСТ~2.105--2019,~--- а также на требования Учёного совета ИТМО к
оформлению ВКР, и по ним настроил параметры документа: формат листа~A4,
поля 30~мм слева и по 20~мм с остальных сторон, шрифт Times New Roman
кеглем~14~пт, полуторный интервал, абзацный отступ 1,25~см, выравнивание
по ширине. Подобрал я и нужные пакеты \LaTeX{}~--- в частности, для
автоматической сборки оглавления и библиографии.

Файлы проекта я разложил по принципу <<один структурный элемент~--- один
файл>>: введение, каждый из четырёх разделов, заключение и приложения
лежат в отдельных файлах, таблицы вынесены в отдельную папку,
иллюстрации~--- в папки со схемами и с графиками отдельно, а настройки
оформления собраны в отдельном заголовочном файле. Главный файл
подключает всё остальное и собирает работу в один PDF. Такая разбивка
сильно ускорила работу: я мог править один раздел и сразу видеть в Git,
что изменился только он, а не весь документ целиком.

Не обошлось без проблем оформления, которые приходилось устранять в
процессе работы. По ходу написания у меня возникали нерегулярные
отступы, растягивание текста в таблицах, неправильные переносы в
специфических местах. В содержании цифры одно- и двузначных номеров
подразделов <<слипались>> с текстом, и мне пришлось задавать им разную
ширину, а в отточиях~--- учащать шаг точек, чтобы оно выглядело по ГОСТу.
Длинные многостраничные таблицы я оформлял так, чтобы шапка
автоматически повторялась на каждой странице, а ширины колонок подгонял
вручную.

\subsection{Согласование темы и задания}

Тему ВКР я согласовал через личный кабинет обучающегося на сайте
\texttt{my.itmo.ru} в разделе <<Диплом~$\to$~Заявление>>. Там я указал
тему работы, вид ВКР~--- классическая, область~--- прикладные
исследования, и руководителя~--- Быковского Сергея Вячеславовича. Само
заявление затем согласовали и подтвердили научный руководитель,
секретарь ГЭК и руководитель образовательной программы.

Структуру работы я сформировал ещё в феврале. Для этого дважды
проводилась очная консультация с научным руководителем с промежутком в
две недели, на которых я задавал вопросы по интересующим меня местам:
какие должны быть разделы, как их лучше расположить, что самое главное
отразить в каждом, как связать разделы между собой, как органично прийти
к заключению, какие взять критерии достижимости цели и на что опираться
при защите.

Параллельно с обсуждениями я занимался непосредственной разработкой
своего решения~--- анализом предметной области и экспериментальной
апробацией. Я исходил из следующего соображения: пока я не получу
практическое подтверждение идеи, не стоит подробно описывать процесс
работы, так как на практике вскрываются подводные камни, из-за которых
текст придётся переписывать. Поэтому я сначала сфокусировался на том,
чтобы довести до работы минимальный прототип и доказать, что цель
достижима, а сам текст вёл параллельно в рамках производственной
практики, в которой эта разработка и шла.

Позже, уже после доработки обзора, я заполнил индивидуальное задание на
ВКР в той же вкладке <<Диплом>> на \texttt{my.itmo.ru}. Я подготовил
варианты формулировок задания в документе формата~\texttt{.docx} и
прислал их научному руководителю; он внёс правки и комментарии, а я
выбрал тот вариант, который лучше всего подходил к моей работе, и
загрузил его на проверку. Задание согласовали научный руководитель,
секретарь ГЭК и руководитель образовательной программы.

\subsection{Сбор и систематизация источников}

После согласования задания я приступил к сбору источников для
текста ВКР. Согласно требованиям индивидуального задания, в
обзоре литературы должно быть использовано не менее
20~источников не старше 5~лет, в том числе не менее~20\% научных
статей в изданиях, рекомендованных ВАК РФ для диссертаций, и
не менее~15\% в изданиях, индексируемых в международных базах
Web of Science и Scopus.

Для проверки наличия журналов в перечне рекомендованных ВАК~РФ
изданий я скачал актуальный перечень с веб-сайта
\texttt{vak.minobrnauki.gov.ru} в формате PDF и сохранил его
локально. Каждое из изданий-кандидатов я искал внутри документа
горячими клавишами Ctrl+F по полному и сокращённому названию
журнала. Для проверки индексации иностранных журналов в Scopus
я использовал сервис \texttt{scopus.com/sources}, в который
вводил название журнала и проверял, что охват источника в
индексе доходит до текущего 2026~года. Дополнительно по каждому
российскому журналу я проверял категорию в перечне ВАК (К1, К2,
К3)~--- более высокая категория указывает на более строгие
требования к публикациям и придаёт ссылке на статью большую
ценность с формальной точки зрения.

Поиск самих источников я вёл несколькими способами параллельно.
Для зарубежных научных статей основными ресурсами служили
\texttt{arxiv.org}, \texttt{ieeexplore.ieee.org},
\texttt{dl.acm.org}, \texttt{scholar.google.com} и
\texttt{proceedings.mlr.press}. По ключевым словам
<<keyword spotting>>, <<embedded inference>>, <<quantization>>,
<<voice activity detection>>, <<DS-CNN>>, <<MFCC>>,
<<TinyML>>, <<MLPerf Tiny>>, <<cascaded detection>>,
<<always-on listening>> я находил статьи в материалах
конференций (CVPR, NeurIPS, MLSys, Interspeech, tinyML
Research Symposium) и в журналах (IEEE Transactions on Circuits
and Systems~I, IEEE Journal of Solid-State Circuits, Future
Internet, EURASIP Journal on Audio, Speech, and Music
Processing). Для российских научных статей в журналах перечня
ВАК~РФ я использовал портал \texttt{elibrary.ru}, выполнял
поиск по ключевым словам <<распознавание ключевых слов>>,
<<квантизация нейронных сетей>>, <<голосовое управление>>,
<<детектор речевой активности>>, <<свёрточные нейронные сети>>,
<<распознавание речи на микроконтроллерах>>, <<встраиваемые
системы>>, <<нейроморфные вычисления>>.

По каждой потенциально подходящей статье я скачивал полный
PDF-документ, читал его и составлял на русском языке короткое
резюме (4--8 предложений) с указанием: какие именно результаты
работы я планирую использовать в тексте ВКР; в каком разделе
эта ссылка появится; на что конкретно я опираюсь~--- метод,
численный результат, формулировку, графическое представление,
методику оценки. Резюме я хранил в виде Markdown-файлов в
структуре, повторяющей содержание ВКР, что обеспечило
прозрачную трассируемость каждой ссылки.

В итоге сформирован список из 29~использованных источников, в
котором:
\begin{itemize}
	\item источников не старше 5~лет (опубликованных в 2021~году
	      и позднее)~--- 21, что превышает минимальное требование
	      <<не менее 20>>;
	\item научных статей в журналах, рекомендованных ВАК~РФ для
	      диссертаций,~--- 7, что составляет 24,1\% от общего
	      числа источников и превышает минимальное требование
	      <<не менее 20\%>>;
	\item статей в изданиях, индексируемых базами Web of
	      Science и Scopus,~--- 13, что составляет 44,8\% от общего
	      числа источников и превышает минимальное требование
	      <<не менее 15\%>>.
\end{itemize}

В составе ВАК-источников представлены публикации в журналах:
<<Научно-технический вестник информационных технологий, механики
и оптики>> (два источника, оба двойной индексации
ВАК~+~Scopus), <<Цифровая обработка сигналов>> (два источника),
<<Автоматика и телемеханика>> (один источник, двойная
индексация), <<Компьютерные исследования и моделирование>>
(один источник, категория~К2, двойная индексация),
<<Компьютерная оптика>> (один источник, категория~К1, двойная
индексация Scopus~Q2). Подобное соотношение обеспечивает
запас выполнения формальных требований по каждой категории
источников на уровне не менее 4\%.

\subsubsection{Ключевые источники, определившие ход работы}

Среди собранных источников выделю те, на которые я опирался при
написании наиболее интенсивно и которые оказали определяющее
влияние на методологию и оформление ВКР.

Наибольшее влияние на работу оказала статья Сёренсена и
др. в журнале EURASIP Journal on Audio, Speech,
and Music Processing, посвящённая развёртыванию разделимой
свёрточной сети для распознавания ключевых слов на встраиваемой
платформе. Работа отличается исключительно наглядным
оформлением результатов, я перечитывал её несколько раз и
заимствовал из неё сразу несколько решений. Во-первых, из неё я
взял метрики и общую методологию систематического исследования
архитектур нейросетевых моделей. Во-вторых, способы графического
представления результатов, в том числе диаграммы размаха (box
plot, в тексте ВКР~--- на русском языке). В-третьих, именно из
этой работы у меня возникло понимание, что для энергетической
декомпозиции достаточно выбрать несколько представительных
архитектур, а не выполнять её сразу для всех обученных моделей.
В-четвёртых, под влиянием этой статьи я разработал послойный
профилировщик модели и провёл анализ задержек на каждом слое,
подтвердив работу SIMD-операций над свёртками и умножениями.

Существенным ориентиром послужил бенчмарк MLPerf Tiny.
Из него я взял понятие энергии одного инференса (мДж на инференс)
как основной метрики оценки модели на целевой платформе,
а также способы сравнения моделей; кроме того, бенчмарк задал
конкретную пороговую границу точности (90\%), которую я
использовал при построении Парето-фронтов в координатах
<<точность~$\times$~размер модели>> и
<<точность~$\times$~энергия инференса>>.

Работа Cerutti и др. о субмилливаттном
распознавании ключевых слов на микроконтроллере оказалась
наиболее трудной для получения~--- я долгое время не мог
скачать её полный текст,~--- однако её основной вывод о
доминирующем вкладе предобработки в энергобюджет KWS-конвейера
стал важной опорой при интерпретации собственных измерений.
Работы Chen и др. и Lin и др. о
сверхнизкопотребляющих аналоговых детекторах голосовой
активности существенно расширили мой кругозор в части предельно
достижимых показателей первой ступени каскада и использованы при
обосновании направлений дальнейшего развития системы.

Отдельно отмечу основополагающую работу Zhang и др.
<<Hello~Edge>>, задавшую эталонную конфигурацию
архитектуры DS-CNN. На неё я опирался при выборе параметров
конвейера предобработки и архитектуры классификатора, а её
эталонную модель отдельно рассматривал в экспериментальной части
как референсную точку: при построении Парето-фронта я показал,
что данная конфигурация не является Парето-оптимальной для моей
задачи на платформе ESP32-S3.

\subsection{Написание разделов работы}

Сам текст я писал по согласованной с руководителем структуре: введение,
четыре раздела основной части с выводами по каждому, заключение, списки
сокращений, терминов и источников, а также приложения.

Во введении я, согласно методическим рекомендациям, дал общую
характеристику работы: описал актуальность темы, степень разработанности
проблемы, конкретизировал решаемую проблему, сформулировал цель и задачи,
а также практическую значимость. По требованиям Учёного совета ИТМО для
бакалаврской ВКР этот перечень элементов является исчерпывающим. После
обсуждения с руководителем я дополнительно добавил во введение два блока,
формально для бакалаврской работы необязательных~--- <<Апробация и
внедрение результатов>> и <<Личный вклад автора>>,~--- поскольку они
делают введение понятнее для комиссии и подчёркивают самостоятельный
характер выполненной работы.

В первом разделе я проводил анализ предметной области: разбирал, как
устроена задача распознавания ключевых слов и какие подходы к её
решению существуют. Именно этот раздел я после второго прослушивания
переписывал заново и расширял (об этом подробнее ниже)~--- добавлял в
него разбор разных классов вычислительных платформ. Второй раздел я
посвятил проектированию своей каскадной архитектуры, третий~---
реализации системы и описанию измерительного стенда, а четвёртый~---
анализу полученных результатов.

Разделы про реализацию и про результаты я переписывал полностью уже на
финальной стадии~--- после того как получил все экспериментальные
данные и согласовал итоговый набор графиков. Раздел с результатами я
переписывал дважды: первый раз~--- по предварительным замерам, а
финальный~--- по полному набору измерений, когда уже можно было
обоснованно выбрать итоговую модель и подвести главный результат
работы. Здесь же сыграли роль требования руководителя к структуре: он
просил, чтобы методология и результаты строго соответствовали друг
другу и шли по перечню задач из введения. Мой первый вариант, где
раздел результатов был свободным изложением, мы отклонили как
недостаточно строгий, и в финальной версии я выстроил оба раздела
ровно по списку задач~--- так комиссии проще проследить, что какой
задаче соответствует.

Заключение я написал по требованиям задания: перечислил основные
результаты по каждой задаче с конкретными числами, обосновал научную
новизну и практическую значимость, отдельно зафиксировал, что цель
достигнута и все задачи выполнены, и в конце расписал направления
дальнейших исследований по теме.

\subsubsection{Введение}

Во введении я дал общую характеристику работы в соответствии с
методическими рекомендациями: описал актуальность темы исследования,
степень разработанности проблемы, решаемую проблему, цель и задачи
работы, а также практическую значимость. Каждый из перечисленных
элементов выделен в отдельный абзац с полужирным заголовком, что
облегчает восприятие структуры введения.

Описывая актуальность, я указал на фундаментальное противоречие
между требованием непрерывного мониторинга аудиопотока (микрофон
должен всегда быть готов к приёму голосовой команды) и
ограниченным энергетическим бюджетом батарейного IoT-устройства
(непрерывная работа аудиотракта приводит к разряду источника
питания в течение единиц часов вместо ожидаемых месяцев и лет).
В качестве альтернативы локальному распознаванию я рассмотрел
облачные голосовые ассистенты (Amazon Alexa, Google Assistant,
Apple Siri) и обосновал их непригодность для батарейных
устройств: высокое потребление радиомодуля при передаче
аудиоданных по беспроводному каналу, сетевая латентность,
зависимость от инфраструктуры, риски для приватности голосовых
данных пользователя.

В блоке <<Степень разработанности>> я систематизировал три
направления повышения энергоэффективности систем распознавания
ключевых слов, представленные в международной литературе:
разработку компактных нейросетевых архитектур, применение
методов квантизации весов и активаций, построение каскадных
схем детекции. Я обратил внимание на то, что в известной
литературе данные направления исследуются преимущественно
изолированно, без раздельной количественной оценки их
совместного вклада в итоговый энергобюджет цикла обработки,
и сформулировал восполнение указанного пробела как
научно-техническую проблему настоящей работы.

Цель работы сформулирована как снижение энергопотребления
устройств интернета вещей с функцией голосового управления за
счёт каскадного управления режимами работы микроконтроллера и
применения квантизованных нейронных сетей для распознавания
ключевых слов голосовых команд. Для достижения цели поставлены
пять задач: проанализировать существующие методы и системы;
разработать каскадную архитектуру детекции; исследовать влияние
методов квантизации на точность, размер и вычислительную
сложность модели; реализовать систему на платформе ESP32-S3 и
провести её экспериментальную апробацию; выполнить декомпозицию
итогового энергобюджета цикла каскадной обработки.

Практическая значимость обоснована применимостью предложенной
архитектуры и методик при проектировании автономных
IoT-устройств с голосовым управлением в задачах промышленного
мониторинга, носимой электроники и беспроводных узлов умного
дома, а также переносимостью архитектуры на микроконтроллеры
семейств STM32, nRF52 и аналогичные с поддержкой аппаратных
режимов глубокого сна и внешнего пробуждения.

После согласования с научным руководителем во введение
дополнительно добавлены два блока, формально не являющиеся
обязательными для бакалаврской ВКР согласно требованиям Учёного
совета ИТМО (где для бакалаврской работы предусмотрены только
актуальность, степень разработанности, решаемая проблема, цель
и задачи, практическая значимость). Это блоки <<Апробация и
внедрение результатов>>~--- с указанием на доведение
предложенной архитектуры до работающего программно-аппаратного
решения, его экспериментальную характеризацию на лабораторном
стенде и размещение исходного кода в открытых удалённых
репозиториях для обеспечения воспроизводимости,~--- и <<Личный
вклад автора>>~--- с фиксацией самостоятельного характера всех
полученных результатов. Оба блока повышают информативность
введения и облегчают комиссии понимание содержания работы ещё
до перехода к основной части.

\subsection{Диаграммы, иллюстрации и листинги}

Отдельно стоит отметить, что создание визуального сопровождения заняло
существенно больше запланированного времени, вследствие чего я загрузил
работу на антиплагиат в последний момент. Диаграммы я готовил с оглядкой
на то, что у нас на кафедре требовали строго соблюдать нотации.
Компонентную диаграмму программного обеспечения я делал в нотации~C4.
Сначала я попробовал рисовать её текстовыми средствами, но
автоматическая раскладка получалась плохо читаемой и трудно правилась
вручную, поэтому финальную диаграмму я собрал вручную во встроенной
библиотеке форм~C4 редактора draw.io. Диаграмму состояний и диаграмму
последовательности, наоборот, я описал текстовой нотацией, поскольку
они ей хорошо ложатся. Принципиальную электрическую схему стенда я
нарисовал в среде KiCad, а графики результатов строил скриптами на
Python и складывал эти скрипты в отдельную папку удалённого репозитория,
чтобы любой график можно было пересобрать из исходных данных. Все
диаграммы я экспортировал в векторном формате и вставлял в \LaTeX{},
чтобы они оставались чёткими при любом масштабе.

Полные листинги кода я не прикладывал к работе ввиду ряда причин.
Исходного кода встроенного программного обеспечения и вспомогательных
скриптов набралось несколько тысяч строк, и приложить их целиком не
представляется возможным с точки зрения объёма; при этом представление
отдельных частей также не является решением проблемы, так как они не
отражают всех особенностей встраиваемого программного обеспечения.
Поэтому я описал структуру программного обеспечения визуально, в
нотации~C4, и словами разобрал ключевые инженерные решения, а в
приложения вынес только самые важные фрагменты кода. Полный код я
выложил в публичные удалённые репозитории и сослался на них.

В приложения я также добавил принципиальную схему стенда, диаграммы
состояний и последовательности, фотографии стенда в разных фазах
работы, таблицу разделов флеш-памяти и ссылки на удалённые репозитории.
Когда вся работа была готова, я написал подробные файлы README.md для
трёх удалённых репозиториев проекта~--- с описанием, что это, из чего
состоит и как собирать,~--- чтобы в проделанной работе можно было
разобраться со стороны.

\subsection{Итерации с руководителем и прослушивания}

Процесс работы над текстом ВКР совместно с научным руководителем можно
разделить на четыре этапа. Первый этап~--- обсуждение структуры в
феврале. Второй~--- разбор содержания после первого прослушивания:
уточнили формулировки цели и задач, поправили обзор. Третий~--- после
второго прослушивания: расширил обзор и переформулировал цель.
Четвёртый~--- финальная проверка работы и согласование задания и
аннотации. Между этими этапами мы переписывались по почте и в
мессенджере по отдельным вопросам.

К первому и второму прослушиванию я готовил презентацию, регламент
выступления был 5~минут. На первом прослушивании (его цель~---
согласовать тему для приказа) надо было показать актуальность темы с
разбором аналогов, ожидаемый практический результат и степень
готовности в процентах, цель и задачи, тему и руководителя. Ко второму
прослушиванию требования были аналогичными, у меня изменился процент
готовности, обновился обзор и список выполненных задач и слегка
изменилось оформление презентации.

На втором прослушивании комиссия задала мне ряд вопросов про
альтернативные вычислительные платформы. Святослав Осипов спросил,
рассматривал ли я ПЛИС и FPGA вместо обычного микроконтроллера. Я
ответил, что специально брал один из самых распространённых
микроконтроллеров ради гибкости и доступности. Он отметил, что слайд с
актуальностью выглядит неполным, поскольку между специализированными
микросхемами и микроконтроллером есть ещё и ПЛИС, и посоветовал либо
подготовиться к таким вопросам, либо прямо в презентации обосновать,
почему ПЛИС не подходит. Василий Пинкевич продолжил тему: кроме ПЛИС
есть ещё цифровые сигнальные процессоры и встроенные графические
ускорители, и для подобных задач они тоже применяются. Я ответил, что
полтора года работал на этом семействе микроконтроллеров и специально
выбрал чип с поддержкой SIMD, а квантизацию весов сделал именно для
того, чтобы пользоваться этими быстрыми операциями. Осипов посоветовал
обосновать выбор через экономику и через тезис, что ПЛИС для такой
лёгкой задачи~--- это <<из пушки по воробьям>>; в остальном работу
оценили хорошо. Василий Пинкевич также спросил, есть ли у меня метрики
качества распознавания кроме энергопотребления, и я показал точность,
F1-меру и матрицу ошибок.

После этого прослушивания я заново полностью прогнал обзор предметной
области в тексте ВКР и добавил в него разбор ПЛИС, DSP, FPGA и других
классов вычислительных платформ~--- по каждому нашёл публикации, разобрал
их ограничения и обосновал, почему для моей задачи подходит именно
микроконтроллер общего назначения. Это расширение я внёс и в раздел
обзора, и в список литературы. Тогда же, в обсуждении задания на ВКР, я
вместе с руководителем поменял формулировку цели и задач.

\subsection{Аннотация и заполнение полей в личном кабинете}

15~мая я заполнил аннотацию к ВКР через личный кабинет
\texttt{my.itmo.ru} в разделе <<Диплом~$\to$~Аннотация>>. В
аннотации я указал цель исследования, перечислил пять задач,
дал краткую характеристику полученных результатов с указанием
ключевых количественных характеристик. После заполнения
аннотация была согласована научным руководителем, секретарём
ГЭК и руководителем образовательной программы.

После согласования аннотации в личном кабинете появилась
возможность загрузить итоговый файл ВКР через вкладку
<<Диплом~$\to$~Файл>>. При нажатии на эту вкладку открывается
форма с полем загрузки файла; после загрузки система
\texttt{my.itmo.ru} автоматически формирует итоговый PDF-документ,
включающий титульный лист, аннотацию, индивидуальное задание
и загруженный текст работы. Эту функцию я использовал на
следующем этапе практики при отправке текста ВКР на
нормоконтроль и проверку в системе <<Антиплагиат>>.

\section{Проверка текста ВКР в системе <<Антиплагиат>>}

\subsection{Подготовка к нормоконтролю}

19~мая я подготовил текст ВКР к отправке на проверку. Поскольку
работа изначально вёрстается в \LaTeX{}, для получения итогового
PDF-файла мне не пришлось выполнять дополнительной конвертации
из других форматов~--- достаточно было полностью пересобрать
проект командой \texttt{latexmk -xelatex main.tex}. Перед
сборкой я ещё раз пересмотрел все разделы на предмет очевидных
опечаток, проверил корректность всех перекрёстных ссылок (на
рисунки, таблицы, формулы и приложения), отсутствие <<висячих>>
строк в конце страниц и корректность переносов в длинных
англоязычных терминах. Итоговый PDF-файл получил имя
\texttt{thesis\_RU.pdf} объёмом~114~страниц.

Сразу после сборки финального PDF-файла я отправил его
секретарю ГЭК Смирновой Анжелике Леонидовне в личные
сообщения мессенджера Telegram с просьбой провести
нормоконтроль перед загрузкой в систему <<Антиплагиат>>.
Контакт секретаря ГЭК я узнал из чата выпускников
образовательной программы в Telegram, где Анжелика Леонидовна
сообщает организационные сведения и отвечает на вопросы
студентов в течение всего семестра.

\subsection{Нормоконтроль и устранение замечаний}

Замечаний по нормоконтролю у меня оказалось мало, поскольку
шаблон \LaTeX{}, на котором я писал работу, изначально
реализует все требования к оформлению. Основное полученное
замечание касалось оформления одной из таблиц в разделе~3, где
у меня в исходном коде \LaTeX{} использовалась команда
\verb|\begin{tabular}| без обёртки в \verb|\begin{longtable}|, и
при автоматическом разрыве таблицы между страницами заголовок
повторялся некорректно. Я заменил всё определение таблицы на
\verb|longtable| с явным указанием заголовков
\verb|\endfirsthead| и \verb|\endhead|, после чего таблица стала
переноситься корректно с автоматическим продолжением и
повторением шапки на каждой странице.

Второе замечание касалось пары мест в тексте, где формулы
располагались сразу после абзаца без вводящей фразы. Согласно
требованиям к оформлению, перед каждой формулой должна быть
вводящая её фраза, а после формулы~--- расшифровка
обозначений, начинающаяся со слова <<где>> со строчной буквы.
Я отредактировал соответствующие фрагменты, после чего
секретарь ГЭК подтвердила, что замечаний больше нет.

\subsection{Загрузка файла и прохождение проверки}

После устранения замечаний по нормоконтролю, 20~мая, я
загрузил итоговый PDF-файл \texttt{thesis\_RU.pdf} в личном
кабинете \texttt{my.itmo.ru} через вкладку
<<Диплом~$\to$~Файл~$\to$~Текст~ВКР>>. После загрузки в той же
панели появилась ссылка <<Скачать итоговый файл ВКР>>, по
которой можно было загрузить сформированный системой
итоговый PDF-документ с титульным листом, аннотацией и
заданием~--- полная версия работы, которая будет передана в
ГЭК.

Секретарь ГЭК в тот же день, 20~мая, провела проверку текста
ВКР в системе <<Антиплагиат.ВУЗ>>. Результаты проверки:
оригинальность~--- 97,36\%, совпадения~--- 2,64\%,
цитирования~--- 0\%, самоцитирования~--- 0\%, ИИ-контент~---
0\%. Требование индивидуального задания~--- оригинальность не
менее 75\% или заимствований не более 25\%~--- выполнено с
запасом более чем на~22\%. Полученная справка
о результатах проверки в системе <<Антиплагиат>> приведена в
приложении~\ref{app:antiplagiat}.

Высокий уровень оригинальности обусловлен самостоятельным
характером проведённой работы~--- оригинальной формулировкой
проектных критериев каскадной архитектуры, самостоятельной
реализацией программного обеспечения микроконтроллера и
лабораторного измерительного стенда, оригинальной методикой
систематического исследования пространства архитектур
DS-CNN со статистической обработкой результатов по
доверительному интервалу Уилсона и авторской декомпозицией
энергобюджета цикла каскадной обработки на платформе ESP32-S3.
Зафиксированные системой совпадения относятся преимущественно
к стандартным техническим формулировкам и определениям
терминов, в которых сложно избежать повторения с известной
литературой (например, описание базовых формул математической
обработки сигналов или стандартные определения понятий
квантизации и каскадной детекции).

\section{Предзащита ВКР}

\subsection{Подготовка электронной презентации}

Презентацию к предзащите я полностью переработал, оставив без изменений
только первый, титульный слайд. Я сделал унифицированный дизайн в
фиолетовом стиле, упростил фон слайдов до простого белого и использовал
заголовок из официального шаблона ИТМО. Нумерацию слайдов я оформил так,
как просил секретарь ГЭК~--- не просто номером слайда, а в формате
<<K/N>> (например, <<5/15>>), то есть какой это слайд и из скольки, чтобы
по ходу выступления было понятно, сколько ещё осталось.

После прошлого прослушивания я понял, что нужно сфокусироваться на
демонстрационной части ВКР: важно, чтобы информация была визуально
репрезентативна и понятна. Поэтому к предзащите я сконцентрировался на
графической репрезентации данных~--- и в самой ВКР, и в презентации,~---
и у меня получилось подготовить множество графиков: круговую диаграмму
для энергобюджета, временную диаграмму тока с декомпозицией по ступеням
каскада, несколько Парето-фронтов, тепловую карту и принципиальную
схему. Сам текст ВКР к этому моменту был готов, я лишь усиливал
визуальную составляющую.

Графики и диаграммы я готовил тем же набором инструментов, что и для
текста ВКР: статистические графики строил на языке Python с библиотеками
matplotlib и seaborn, диаграммы состояний и последовательности~---
текстовой нотацией Mermaid, принципиальную электрическую схему стенда~---
в среде KiCad, а компонентные и концептуальные диаграммы~--- в редакторе
draw.io. Скрипты построения графиков я держал в отдельной директории
удалённого репозитория ВКР, чтобы каждый график можно было пересобрать из
исходных CSV-данных измерений.

Саму презентацию я оформлял, как и презентации для прослушиваний, в
Google~Slides. Шаблоны брал официальные, от Университета ИТМО. В
Google~Slides я делаю все презентации последние пять лет, так как это
наиболее удобный для меня инструмент с отличной переносимостью: он гибок
и удобен, в отличие от тяжёлого формата \texttt{.pptx}, для которого
часто нужен проприетарный софт и в котором нет такой удобной поддержки
совместной работы~--- возможности показать презентацию другим и спросить
совета по оформлению слайдов. Для справки: шрифт Golos~Text является
предустановленным в шаблоне презентации, и я не стал его менять, так как
он показался мне наиболее визуально приятным из шрифтов без засечек. По
наполнению я старался действовать, опираясь на максимальную прозрачность,
наглядность и простоту~--- задача была показать, что я достиг цели по
конкретным достижимым критериям и выполнил все задачи, и это я во всей
презентации и демонстрировал за счёт графиков и схем.

\subsubsection{Подготовка демонстрационного видеоролика}

Отдельным и наиболее трудоёмким элементом подготовки презентации стала
запись демонстрационного видеоролика работы устройства. Подготовку я
начал с написания сценария: я взял тетрадь и ручку и составлял сценарий
для короткого ролика длительностью около 30~секунд. Я прокручивал в
голове, что именно наиболее релевантно показать комиссии~--- какие
локации, какую функциональность устройства, какие его части. По итогам
этого анализа я сформировал ряд эвристик для ролика. Основной вывод
состоял в том, что в текущей реализации устройства недостаточно
визуальной составляющей, по которой можно было бы наглядно проследить
его работу.

Чтобы устранить этот недостаток, я взял имеющиеся у меня IPS-дисплеи
разрешением 240$\times$240, подключил их по интерфейсу SPI и написал
программное обеспечение для вывода шрифтов, после чего синхронизировал
состояния доменов измерительного стенда с дисплеями. На первый дисплей
(домен измеряемой системы) выводились состояния конечного автомата FSM с
интерактивной индикацией загрузки, цветовой кодировкой состояний,
информацией о времени прохождения каждого этапа и уверенности модели в
классификации после инференса. На второй дисплей (домен измерителя
энергопотребления) я в реальном времени выводил график мгновенного тока
$I(t)$ измеряемого домена в диапазоне 0--100~мА, что позволяло
непосредственно наблюдать потребление устройства в процессе работы.

Записанный первоначальный ролик я отправил научному руководителю, от
которого получил замечание: ролик следовало сократить до 15~секунд и
показать лишь самое важное. На основании этого я заново проработал
сценарий и пришёл к выводу, что достаточно показать включение
лабораторного стенда и две успешные классификации. После повторной
записи я намеренно ускорил моменты простоя~--- например, запись голоса и
циклы засыпания и пробуждения устройства,~--- что отдельно пояснил
комиссии в тексте выступления. Дополнительно я добавил в ролик текст с
транскрипцией речи на случай, если на презентации не будет звука или
кто-то не расслышит сказанное.

Монтаж ролика я выполнил в видеоредакторе YouCut на мобильном телефоне
Samsung~A52, выставил максимальную громкость и высокое качество видео,
после чего загрузил его на Google~Диск и добавил в презентацию в
Google~Slides как видео. Я отдельно убедился, что ролик действительно
воспроизводится и что его хорошо слышно и видно.

\subsubsection{Содержание слайдов презентации}

Презентация состоит из 15~содержательных слайдов и завершающего слайда
<<Спасибо за внимание>> (всего 16). Ниже описаны основные решения по
каждому слайду; сами слайды приведены в
приложении~\ref{app:presentation}.

\textit{Слайд~1 (титульный).} Титульный слайд я перенёс без изменений со
второго прослушивания. Тема ВКР с того времени не менялась, поскольку
мы с научным руководителем согласовали удачную формулировку. Изменения
вносились лишь по итогам первого прослушивания: термин <<IoT>> был
заменён на <<устройства интернета вещей>>, а <<распознавание речи>>~---
на <<распознавание голосовых команд>>, что точнее описывает суть работы,
поскольку распознавание произвольной речи~--- существенно более сложная
задача, тогда как я распознаю ключевые слова голосовых команд.

\textit{Слайд~2 (актуальность и проблематика).} Над этим слайдом я
работал на каждом прослушивании. После первого прослушивания мне
указали на недостаточный обзор электроники, после второго~--- на нехватку
разбора DSP и FPGA и слабую защиту по этим пунктам. К предзащите я
рассмотрел эти классы платформ на разных уровнях, однако основную их
часть решил перенести в слайд, раскрывающий первую задачу (анализ
предметной области), а слайд об актуальности оставить более простым и
сфокусировать его на проблеме автономности батарейных устройств. На
слайде я показал, что голосовое управление становится стандартом для
целого класса устройств интернета вещей~--- носимой электроники, умного
дома, промышленных датчиков, медицинских приборов,~--- которые требуют
автономности и работают от батареи, вследствие чего распознавание команд
должно происходить локально, без отправки звука в облако. Локальное
распознавание означает непрерывное прослушивание звукового потока, что
требует постоянной работы микрофона и процессора и вступает в прямой
конфликт с ограниченным энергетическим бюджетом батареи~--- это и есть
проблема, решаемая в работе.

\textit{Слайд~3 (цель и критерии достижимости).} Цель работы я поместил
в выделенную фиолетовую рамку для акцентирования внимания, а критерии
достижимости цели перечислил блоками~C1--C5 с подсветкой их конкретных
численных значений.

\textit{Слайд~4 (задачи).} На этом слайде перечислены пять задач,
согласованных в задании на ВКР. Слайд оформлен без дополнительных
визуальных решений: я проверил корректность шрифтов, при выступлении
указывал комиссии, что задачи представлены на слайде, и переходил далее.

\textit{Слайд~5 (анализ предметной области, задача~1).} Это один из
наиболее объёмных слайдов. На нём я отразил, что было проанализировано,
что выявлено и что сформулировано в рамках первой задачи. Сложность
состояла в том, чтобы выразить за 30--40~секунд полноценный анализ, на
который в ВКР ушло несколько месяцев и порядка 10--15~страниц текста.
Структура слайда отражает три выполненных действия. Первое~---
систематизация подходов по двум уровням: по месту выполнения
распознавания (облако против устройства) и по типу вычислительной
платформы (специализированные микросхемы, ПЛИС и сигнальные процессоры,
микроконтроллеры общего назначения). Я обосновал, что облачный подход
непригоден из-за высокого потребления радиомодуля; что специализированные
микросхемы дают рекордную энергоэффективность, но менее гибки в выборе
архитектуры нейросети; что ПЛИС и сигнальные процессоры требуют
разработки аппаратной логики на специализированных языках, удлиняющей
каждую итерацию модели; и что микроконтроллеры общего назначения дают
полную гибкость и быстрые итерации, а моя работа сокращает их разрыв в
энергоэффективности. Второе действие~--- выявление четырёх источников
избыточного потребления: непрерывный захват аудио, постоянная работа
процессора, избыточная сложность моделей в литературе и неиспользование
режимов глубокого сна. Третье~--- формулировка четырёх требований к
архитектуре.

\textit{Слайд~6 (каскадная архитектура).} На этом слайде я
сконцентрировался на максимально наглядном представлении каждой ступени
каскада, сопроводив каждую ступень иллюстрацией происходящего на ней. Для
первой ступени (аналоговая детекция) я показал график зависимости
амплитуды от времени с сигналом команды <<yes>>; для второй (извлечение
признаков MFCC)~--- диаграмму <<мел-канал~$\times$~кадр>>; для третьей
(классификация)~--- изображение фильтров DS-CNN и финального слоя
softmax с выбором команды <<yes>>.

\textit{Слайд~7 (исследование архитектур).} На этом слайде представлен
Парето-фронт зависимости точности классификации от размера INT8-модели~---
один из наиболее трудоёмких графиков работы, для построения которого я
обучил около 150~архитектур и выполнил столько же квантизаций. Задача
состояла в том, чтобы построить Парето-фронт, сделать вывод о модели из
работы Hello~Edge, выбрать собственную оптимальную по совокупности
архитектуры, размера и точности модель, а также из интереса выбрать более
компактную модель из области ярко выраженного <<колена>> (knee point)
Парето-фронта. Справа я разместил тепловую карту точности в пространстве
числа блоков и фильтров, показывающую, что точность существенно
деградирует при числе фильтров ниже примерно 48; что даже при двух блоках
точность может быть высокой при 224~фильтрах; и что при большом числе
блоков (например, восьми) всё равно требуется более 40~фильтров для
достижения точности хотя бы в 94\%. Наиболее оптимальные варианты
сгруппировались на границе пяти блоков и более чем 96~фильтров, а наиболее
компактные~--- в области 3--4~блоков и порядка 80--88~фильтров.

\textit{Слайд~8 (сравнение методов квантизации).} На этом слайде я
наглядно показал разницу точности между квантизацией QAT и PTQ и
доверительный интервал, в который попадает большинство моделей, и сделал
вывод, что QAT не даёт существенного преимущества над PTQ. Для
наглядности я выбрал несколько моделей различных размеров и категорий и
для каждой показал три столбца~--- FP32, INT8~PTQ, INT8~QAT~--- с
доверительным интервалом. Получилось два графика друг над другом, а
справа~--- краткие выводы. В заметках к слайду я записал текст
выступления, чтобы не сбиваться и не повторяться.

\textit{Слайд~9 (выбор финальной модели).} На этом слайде по оси абсцисс
отложена энергия одного инференса в мДж~--- ключевая для работы величина,
заимствованная из изученных статей и бенчмарков. Это один из главных
выводов работы: на Парето-фронте выбрана итоговая модель с 104~фильтрами
и 5~блоками, оказавшаяся наиболее экономичной, тогда как модель с
176~фильтрами и 7~блоками потребляет втрое больше при разнице в точности
всего 0,17\%. Выбранная модель располагается в характерной точке
<<колена>> Парето-фронта.

\textit{Слайд~10 (лабораторный измерительный стенд).} На этом слайде я
показал подготовленный демонстрационный видеоролик и принципиальную схему
стенда, составленную в KiCad. Я стремился оформить слайд максимально
наглядно, чтобы продемонстрировать инженерные навыки, однако на
предзащите получил замечания: рамка принципиальной схемы выполнена не по
ГОСТ и сама схема сложна для прочтения. Поэтому впоследствии я полностью
переработал схему; исправленный вариант слайда приведён в
приложении~\ref{app:fixed_slide}. (На предзащите слайд назывался <<слайд
с двухдоменным измерительным стендом>>; в исправленной версии я
переименовал его в <<лабораторный стенд>>.)

\textit{Слайд~11 (профиль тока каскада).} На этом слайде размещены два
графика друг над другом во весь слайд: сверху~--- исходные (необработанные)
данные $I(t)$ с монитора питания, снизу~--- усреднённые по фазам значения,
показывающие конкретное потребление на каждой фазе конечного автомата.

\textit{Слайд~12 (декомпозиция энергобюджета).} Это один из важнейших
слайдов. Для него я подготовил множество вариантов круговых диаграмм и в
итоге остановился на разноцветной диаграмме с текстовыми выносками для
каждой составляющей. Задача состояла в том, чтобы показать вклад каждой
фазы работы микроконтроллера. Предварительно я выполнил обработку и
исключил из диаграммы этапы индикации результата и глубокого сна,
поскольку их вклад невозможно оценить детерминированно: индикация нужна
сугубо для демонстрационного ролика, а длительность глубокого сна
варьируется в зависимости от сценария использования. Основная мысль~---
показать этапы, которые присутствуют в каждом цикле работы каскада в
неизменном виде, и выявить, какие из них вносят наибольший вклад в
потребление и потому подлежат оптимизации в первую очередь. Получилось,
что этап записи звука длительностью 2,5~с расходует около 569~мДж, а этап
пробуждения занимает около двух секунд и расходует порядка 1982~мДж, что
непозволительно много.

\textit{Слайд~13 (автономная работа от аккумулятора).} На этом слайде я
оценил автономность реализованного устройства при промышленной
эксплуатации. За основу я взял аккумулятор ёмкостью 2500~мА$\cdot$ч,
поскольку именно такая батарея использовалась на стенде, а аккумуляторы
формата~18650 близкой ёмкости широко распространены. По оси ординат
отложена автономность в сутках в логарифмическом масштабе (степени
десяти), столбцы построены для стендовой отладочной платы, для целевой
системы и для теоретического минимума на кристалле в режиме RTC-only.
Для целевой системы с ULP~VAD автономность при одном пробуждении в час
составляет около двух лет~--- основной демонстрируемый результат:
грамотное сочетание каскада и режимов управления питанием позволяет
устройству засыпать, ожидать пробуждения и работать годами от обычного
аккумулятора.

\textit{Слайд~14 (достижение цели и новизна).} На этом слайде я привёл
критерии достижимости цели, показал достижение цели по каждому критерию
и тезисно сформулировал научную новизну работы~--- декомпозиция
энергобюджета, систематическое исследование архитектур и сравнение
методов квантизации.

\textit{Слайд~15 (перспективы).} Содержание этого слайда я обдумывал
дольше всего. В качестве основных направлений я вынес: сокращение этапа
инициализации, который, как показала декомпозиция, занимает
неоправданно много времени (соответствующие подходы реализуются в
исследовательских работах с нейросетями на конденсаторах, RTC~SRAM и
т.п.); промышленную реализацию платы с экономичными LDO и MEMS-сенсором
пробуждения по звуку (wake-on-sound) вместо неэффективного LM393;
возможность дистилляции знаний; а также необходимость верификации
подхода на других платформах (STM32U5, nRF52840), несмотря на то что на
целевой платформе решение уже работает.

\textit{Слайд~16 (спасибо за внимание).} Завершающий слайд я перенёс с
прослушиваний; на нём указаны мой Telegram и адрес электронной почты для
связи.

\subsection{Прохождение предзащиты}

Предзащита проводилась 21~мая в дистанционном формате через
видеоконференцию. За полчаса до начала я подключился по
присланной в чате выпускников ссылке, проверил работу
микрофона и демонстрации экрана. Выступление прошло в штатном
режиме~--- я последовательно прошёл по пятнадцати слайдам
основной части, акцентируя внимание комиссии на ключевых
числовых результатах работы (точность~96,12\%, размер
модели~101~КБ, доминирующий вклад сервисных этапов на
уровне~80,8\%, расчётная автономность от единиц суток до
двух лет в зависимости от варианта реализации). Уложился
примерно в семь минут.

Результаты предзащиты:
\begin{itemize}
	\item готовность ВКР оценена в~95\%;
	\item готовность презентации оценена в~85\%;
	\item допуск к защите получен на уровне~80\%;
	\item явка на предзащиту~--- присутствовал;
	\item оригинальность работы по системе <<Антиплагиат>>~---
	      97,36\% (норма~$\geq~75\%$);
\end{itemize}

В целом комиссия отнеслась к работе положительно, отметила
качество и системность проведённого исследования, наличие
большой экспериментальной базы (134~архитектуры, 239~точек),
а также практическую ценность разработанного измерительного
стенда. Решение о допуске к защите принято положительно.
Единственное замечание касалось презентации~--- комиссия
предложила доработать отдельные слайды.

\subsection{Замечания научного руководителя по итогам предзащиты}

После завершения предзащиты, в тот же вечер, научный
руководитель прислал отдельным сообщением подробные
комментарии по итогам моего выступления. Основные замечания
касались презентации и состояли в следующем.

\textit{По слайду с задачами.} Слайд с пятью задачами в виде
длинного нумерованного списка получился перегруженным мелким
текстом, поскольку каждая задача занимает 2--3 строки. На
защите этот слайд читать членам комиссии будет сложно. Решение:
сократить формулировки задач до короткой ключевой фразы
(2--4 слова) и сопроводить каждую задачу пиктограммой,
визуализирующей её содержание (увеличительное стекло для
анализа, схема каскада для архитектуры, шкала для
квантизации, паяльник для реализации, диаграмма для
декомпозиции). Полные формулировки задач сохранить в
заметках к слайду на случай уточняющих вопросов.

\textit{По слайду с Парето-границей <<точность $\times$
	размер>>.} Тепловая карта в правой части слайда оказалась
слишком детальной для презентации~--- 21~строка по числу
фильтров умноженная на~7~столбцов по числу блоков даёт~147
ячеек, разглядеть значения в которых на удалении невозможно.
Решение: оставить только Парето-границу в координатах
<<точность $\times$ размер>>, а тепловую карту перенести в
слайды-приложения для ответов на вопросы. На самой
Парето-границе явно подсветить три ключевые точки~---
выбранную финальную модель $f104\_b5$, эталонную Hello~Edge
$f176\_b6$ и порог MLPerf~Tiny~90\%.

\textit{По слайду с декомпозицией энергобюджета.} Круговая
диаграмма получилась информативной, но подписи долей мелкие
и плохо читаются. Решение: увеличить размер шрифта подписей,
вынести их на внешние выноски с указанием абсолютной энергии
этапа в~мДж и относительной доли в~процентах. Подсвеченный
блок <<новизна>> внизу слайда оставить, но переформулировать
короче и сильнее~--- акцент на том, что именно идентификация
сервисных этапов как доминирующих составляет вклад работы.

\textit{По слайду с автономностью.} Логарифмическая шкала на
вертикальной оси правильная по сути, но без выраженного
обозначения легенды и без подписей числовых значений над
столбцами комиссия может неправильно интерпретировать
визуальный масштаб. Решение: над каждым столбцом подписать
численное значение в человекочитаемом виде (<<17~суток>>,
<<2,9~месяцев>>, <<1,9~года>>, <<2,3~года>>) и явно
обозначить, что вертикальная шкала логарифмическая.

\textit{По завершающим слайдам.} На слайде <<Достижение цели
и новизна>> зелёная плашка <<ЦЕЛЬ ДОСТИГНУТА>> внизу
визуально перекрывает номер страницы. Решение: сдвинуть
плашку выше на несколько миллиметров. На слайде
<<Перспективы>> бейдж <<приоритет~№1>> на блоке про
сокращение инициализации хорошо привлекает внимание, но при
этом другие три направления развития выглядят равнозначными,
хотя на самом деле имеют разную степень проработки. Решение:
у двух из трёх оставшихся направлений добавить мелкие
пометки <<средний приоритет>> и <<долгосрочно>>, чтобы
расставить акценты.

\textit{По общей структуре доклада.} Научный руководитель
отметил, что в моём выступлении хорошо была раскрыта
техническая часть (архитектура, эксперимент, декомпозиция),
но в начале не хватало одной сильной фразы, которая бы
объяснила, в чём конкретно новизна работы для комиссии,
которая может не быть глубоко погружена в тему встраиваемых
нейросетевых систем. Решение: добавить в начале выступления
(после слайда с актуальностью, до перехода к задачам)
одно-два предложения, которые бы простым языком сформулировали
суть исследовательского вклада: <<Существующие работы
рассматривают каскадную архитектуру, квантизацию и
оптимизацию модели изолированно. В моей работе впервые
проведена количественная оценка их совместного применения
на единой платформе, и впервые показано, что сервисные
этапы~--- инициализация и захват аудио~--- составляют 80\%
энергобюджета каскадного цикла, а не вычисления, как
обычно предполагается>>.

\subsection{План доработки презентации}

На основании полученных от научного руководителя комментариев
я составил план доработки презентации к итоговой защите,
которая состоится в начале июня. План включает:

\begin{enumerate}
	\item упрощение слайда с Парето-границей по размеру~---
	      перенос тепловой карты в слайды-приложения, подсветка
	      выводов на оставшейся Парето-границе;
	\item добавление в выступление вступительной фразы о
	      сути исследовательского вклада простым языком после
	      слайда с актуальностью;
	\item проведение полных репетиций выступления с
	      хронометражем для соответствия временным ограничениям
	      процедуры защиты ВКР.
\end{enumerate}

Доработка презентации ведётся параллельно с подготовкой
настоящего отчёта по преддипломной практике и должна быть
завершена не позднее чем за неделю до даты защиты, чтобы
осталось время на финальные репетиции выступления.

Презентация для предзащиты в её исходном виде, которая
обсуждалась на процедуре предзащиты и по которой получены
описанные выше замечания, приведена в
приложении~\ref{app:presentation}.
